% Topic 7: Using Matrix Algebra for Systems of Linear Equations in More Than Two Variables
% Part of the Matrix Fundamentals section

\subsection{Using Matrix Algebra for Systems of Linear Equations in More Than Two Variables}

The elegance of matrix algebra is that the fundamental equation used to solve a system, $X = A^{-1}B$, remains exactly the same regardless of whether you are solving a $2 \times 2$ system, a $3 \times 3$ system, or a $100 \times 100$ system. 

The only difference is the computational effort required to find the inverse matrix, $A^{-1}$.

\bigskip

% ── Worked Example: Solving a 3x3 System using Matrices ──
\begin{problem}[Matrix Solution to a $3 \times 3$ System]
Recall the bakery word problem from the previous topic, which gave us the following system:
\begin{align*}
2x + 3y + z &= 115 \\
x + 2y + 2z &= 110 \\
3x + y + 3z &= 165
\end{align*}
Use matrix algebra to solve for $x, y$, and $z$, finding the exact price of each cake.
\end{problem}

\begin{solution}
\textbf{Step 1: Write the system in the form $AX = B$.} \\
\[
\begin{bmatrix} 2 & 3 & 1 \\ 1 & 2 & 2 \\ 3 & 1 & 3 \end{bmatrix} \begin{bmatrix} x \\ y \\ z \end{bmatrix} = \begin{bmatrix} 115 \\ 110 \\ 165 \end{bmatrix}
\]
Where $A$ is the $3 \times 3$ coefficient matrix, and $B$ is the constants matrix.

\medskip
\textbf{Step 2: Find the determinant of $A$.} \\
Expanding along the first row:
\begin{align*}
\det(A) &= 2 \begin{vmatrix} 2 & 2 \\ 1 & 3 \end{vmatrix} - 3 \begin{vmatrix} 1 & 2 \\ 3 & 3 \end{vmatrix} + 1 \begin{vmatrix} 1 & 2 \\ 3 & 1 \end{vmatrix} \\[6pt]
&= 2(6 - 2) - 3(3 - 6) + 1(1 - 6) \\
&= 2(4) - 3(-3) + 1(-5) \\
&= 8 + 9 - 5 = 12
\end{align*}
Since $\det(A) = 12 \neq 0$, a unique solution exists.

\medskip
\textbf{Step 3: Find the matrix of cofactors, $C$.} \\
Calculating the 9 cofactors (with the checkerboard sign pattern):
\begin{align*}
C_{11} &= +(6 - 2) = 4,   & C_{12} &= -(3 - 6) = 3,  & C_{13} &= +(1 - 6) = -5 \\
C_{21} &= -(9 - 1) = -8,  & C_{22} &= +(6 - 3) = 3,  & C_{23} &= -(2 - 9) = 7 \\
C_{31} &= +(6 - 2) = 4,   & C_{32} &= -(4 - 1) = -3, & C_{33} &= +(4 - 3) = 1
\end{align*}
So, the cofactor matrix is $C = \begin{bmatrix} 4 & 3 & -5 \\ -8 & 3 & 7 \\ 4 & -3 & 1 \end{bmatrix}$.

\medskip
\textbf{Step 4: Find the adjugate and inverse of $A$.} \\
Transpose $C$ to get the adjugate:
\[
\text{adj}(A) = C^T = \begin{bmatrix} 4 & -8 & 4 \\ 3 & 3 & -3 \\ -5 & 7 & 1 \end{bmatrix}
\]
The inverse is $A^{-1} = \frac{1}{\det(A)}\text{adj}(A) = \frac{1}{12} \begin{bmatrix} 4 & -8 & 4 \\ 3 & 3 & -3 \\ -5 & 7 & 1 \end{bmatrix}$.

\medskip
\textbf{Step 5: Solve for $X$ using $X = A^{-1}B$.}
\[
X = \frac{1}{12} \begin{bmatrix} 4 & -8 & 4 \\ 3 & 3 & -3 \\ -5 & 7 & 1 \end{bmatrix} \begin{bmatrix} 115 \\ 110 \\ 165 \end{bmatrix}
\]
Calculate the matrix product first:
\begin{align*}
X &= \frac{1}{12} \begin{bmatrix} 4(115) + (-8)(110) + 4(165) \\ 3(115) + 3(110) + (-3)(165) \\ -5(115) + 7(110) + 1(165) \end{bmatrix} \\[6pt]
X &= \frac{1}{12} \begin{bmatrix} 460 - 880 + 660 \\ 345 + 330 - 495 \\ -575 + 770 + 165 \end{bmatrix} = \frac{1}{12} \begin{bmatrix} 240 \\ 180 \\ 360 \end{bmatrix} = \begin{bmatrix} 20 \\ 15 \\ 30 \end{bmatrix}
\end{align*}

\medskip
\textbf{Conclusion:} \\
The variables are $x=20, y=15, z=30$. \\
The prices are \$20 for a chocolate cake, \$15 for a vanilla cake, and \$30 for a strawberry cake.
\end{solution}

\begin{takeaways}
\item \textbf{Universal Logic:} The framework $X = A^{-1}B$ scales to any number of dimensions. Only the mechanics of finding $A^{-1}$ change.
\item \textbf{Computational Load:} As systems get larger (e.g., $4 \times 4$), finding inverses by hand becomes highly prone to arithmetic errors. This is why computers and graphics cards are used to solve these systems in real life!
\item \textbf{Checking your work:} You can easily verify the final solution by substituting the values back into any of the original equations (e.g., $2(20) + 3(15) + 30 = 40 + 45 + 30 = 115$, which matches Monday's total).
\end{takeaways}
